\begin{table*}[t] 
        \centering
        
        \setlength{\tabcolsep}{8pt}
        \begin{tabular}{l|c|c|cc|cccc}
        \toprule[1.5pt]
        Method & Sensors & WM-Based & NC & DAC & TTC & Comf. & EP & PDMS \\
        \midrule
        TransFuser~\cite{transfuser} & MV \& L & $\times$ & 97.7 & 92.8 & 92.8 & 100.0 & 79.2 & 84.0 \\
        Hydra-MDP~\cite{hydra} & MV \& L & $\times$ & 98.3 & 96.0 & 94.6 & 100.0 & 78.7 & 86.5 \\
        DiffusionDrive~\cite{diffusiondrive} & MV \& L & $\times$ & 98.2 & 96.2 & 94.7 & 100.0 & 82.2 & 88.1 \\
        WoTE~\cite{li2025wote} & MV \& L & $\checkmark$ & 98.5 & 96.8 & 94.9 & 99.9 & 81.9 & 88.3 \\
        \midrule
        UniAD~\cite{uniad} & MV & $\times$ & 97.8 & 91.9 & 92.9 & 100.0 & 78.8 & 83.4 \\
        PARA-Drive~\cite{paradrive} & MV & $\times$ & 97.9 & 92.4 & 93.0 & 99.8 & 79.3 & 84.0 \\
        LAW~\cite{law} & MV & $\checkmark$ & 96.4 & 95.4 & 88.7 & 99.9 & 81.7 & 84.6 \\
        World4Drive~\cite{world4drive} & MV & $\checkmark$& 97.4 & 94.3 & 92.8 & 100.0 & 79.9 & 85.1 \\
        FSDrive~\cite{zeng2025fsdrive} & MV & $\checkmark$& 98.2 & 93.8 & 93.3 & 99.9 & 80.1 & 85.1 \\
        DrivingGPT~\cite{drivinggpt} & SV & $\checkmark$ & 98.9 & 90.7 & 94.9 & 95.6 & 79.7 & 82.4 \\
        Epona~\cite{epona} & SV & $\checkmark$ & 97.9 & 95.1 & 93.8 & 99.9 & 80.4 & 86.2 \\
        ImagiDrive~\cite{li2025imagidrive} & SV & $\checkmark$ & 98.6 & 96.2 & 94.5 & 100.0 & 80.5 & 87.4 \\
        \rowcolor{gray!20}
        WorldDrive & SV & $\checkmark$ & 98.4 & 96.2 & 95.1  & 100.0 & 81.9 & \textbf{88.1} \\
        \midrule
        PWM~\cite{zhao2025pwm}$\dagger$ & SV & $\checkmark$ & 98.6 & 95.9 & 95.4 & 100.0 & 81.8 & 88.1 \\
        DriveVLA-W0~\cite{drivevlaw0}$\dagger$ & SV & $\checkmark$ & 98.7 & 96.2 & 95.5 & 100.0 & 82.2 & 88.4 \\
        \rowcolor{gray!20}
        WorldDrive$\dagger$ & SV & $\checkmark$ & 98.4 & 96.8 & 95.2  & 100.0 & 83.3 & \textbf{89.0} \\
        \midrule
        DriveVLA-W0$\ddagger$ & SV & $\checkmark$ & 99.3 & 97.4 & 97.0 & 100.0 & 88.3 & 93.0 \\
        \rowcolor{gray!20}
        WorldDrive$\ddagger$ & SV & $\checkmark$ & 99.3 & 98.8 & 97.9 & 100.0 & 88.3 & \textbf{93.6} \\
        \bottomrule[1.5pt]
        \end{tabular}
\caption{\textbf{NAVSIM navtest split comparison}.~PDMS and sub-scores reflecting closed-loop performance. MV: multi-view cameras; SV: single-view camera; L: LiDAR. $\dagger$: Training with full navtrain split. $\ddagger$: Best-of-6 performance with oracle.}
\label{tab:closeloop_turn}
\end{table*}

\begin{table*}[h!]
\centering
\setlength{\tabcolsep}{6pt}
\resizebox{0.95\textwidth}{!}{\begin{tabular}{l| c|c | c c c c c c c c c |c}
    \toprule[1.5pt]
    Method 
    & Sensors
    & Stage
    & $\text{NC}$
    & $\text{DAC}$
    & $\text{DDC}$
    & $\text{TLC}$
    & $\text{EP}$
    & $\text{TTC}$
    & $\text{LK}$ 
    & $\text{HC}$ 
    & $\text{EC}$ 
    & $\text{EPDMS}$  \\
    \midrule
    
    LTF\cite{transfuser} & MV &   \makecell{S1 \\ S2} &
    \makecell{96.2 \\ 77.7} & \makecell{79.5 \\ 70.2} & \makecell{99.1 \\ 84.2} & \makecell{99.5\\ 98.0} & \makecell{84.1 \\ 85.1} & \makecell{95.1 \\ 75.6} & \makecell{94.2 \\ 45.4} & \makecell{97.5 \\ 95.7} & \makecell{79.1 \\ 75.9} & 23.1    \\
    \midrule
    DiffusionDrive\cite{diffusiondrive}  & MV &   \makecell{S1 \\ S2} & \makecell{96.8 \\ 80.1} & \makecell{86.0 \\ 72.8} & \makecell{98.8 \\ 84.4} & \makecell{99.3 \\ 98.4} & \makecell{84.0 \\ 85.9} & \makecell{95.8 \\ 76.6} & \makecell{96.7 \\ 46.4} & \makecell{97.6 \\ 96.3} & \makecell{79.6 \\ 72.8} & 27.5    \\
    \midrule
    WorldDrive  & SV &   \makecell{S1 \\ S2} &
    \makecell{97.3 \\ 91.4} & \makecell{89.1 \\ 82.0} & \makecell{97.6 \\ 91.0} & \makecell{99.7 \\ 98.5} & \makecell{60.5 \\ 53.1} & \makecell{96.8 \\ 90.6} & \makecell{87.7 \\ 52.3} & \makecell{93.1 \\ 93.3} & \makecell{60.0 \\ 62.8} & \textbf{34.9}  \\
    \bottomrule[1.5pt]
\end{tabular}}
\caption{\textbf{NAVSIM-v2 navhard split comparison}. EPDMS and sub-scores reflecting closed-loop performance. MV: multi-view cameras; SV: single-view camera.}
\label{table:navhar}
\end{table*}

\begin{table}[ht] 
        \centering
        
        \setlength{\tabcolsep}{6pt}
        \resizebox{0.47\textwidth}{!}{
        \begin{tabular}{l|c|cc|cc}
        \toprule[1.5pt]
        \multirow{2}{*}{Method} &
        \multirow{2}{*}{Sensors} & 
        \multicolumn{2}{c|}{L2~($m$) $\downarrow$} & 
        \multicolumn{2}{c}{CR~(\%) $\downarrow$} \\
         &  & 3$s$ & Avg.  & 3$s$ & Avg.\\
        \midrule 
         UniAD~\cite{uniad} & V &  1.04  & 0.73 &  0.63 & 0.61 \\
         VAD~\cite{vad} & V  & 1.05 & 0.72  & 0.43 & 0.21\\
         SparseDrive~\cite{sparsedrive} & V  & 0.96 & 0.61  & 0.18 & 0.08\\
         \midrule
         LAW~\cite{law} & V  & 1.01 & 0.61  & 0.54 & 0.30 \\
         World4Drive~\cite{world4drive} & V  & 0.81 & 0.50  & 0.33 & 0.16 \\
         Drive-WM~\cite{drivewm} & V  & 1.20 & 0.80  & 0.48 & 0.26 \\
         Epona~\cite{epona} & SV  & 1.98 & 1.25  & 0.85 & 0.36 \\
         WorldDrive & SV  & 0.68 & 0.42  & 0.38 & 0.16 \\
         \bottomrule[1.5pt]
        \end{tabular}}
\caption{\textbf{nuScenes validation split open-loop planning comparison}.~We follow the SparseDrive evaluation metric. V: multi-view camera; SV: single-view camera}
\label{tab:openloop_turn}
\end{table}

\section{Experiments}


\subsection{Evaluation on Trajectory Planning}
For the end-to-end planning task, we use the NAVSIM~\cite{navsim}, NAVSIM-v2~\cite{cao2025pseudo} and nuScenes~\cite{nuscenes} benchmarks to evaluate the planning performance. NAVSIM utilizes the Predictive Driver Model Score (PDMS) as the closed-loop planning metric, which includes five sub-scores: no-at-fault collisions (NC), drivable area compliance (DAC), time-to-collision (TTC), comfort (Comf.), and ego progress (EP). NAVSIM-v2 further introduces reactive traffic with pseudo closed-loop simulation, and extends PDMS with additional compliance and comfort metrics, including traffic light compliance (TLC), driving direction compliance (DDC), lane keeping (LK), history comfort (HC), and extended comfort (EC). Together, these metrics provide a more comprehensive assessment of closed-loop driving behavior. We also utilize the nuScenes benchmark for a supplementary comparison to evaluate the performance of WorldDrive. We employ the L2 distance error and the collision rate as metrics to assess the performance of open-loop planning.

\subsection{Evaluation on Scene Generation}
To assess the generation quality of the TA-DWM, we evaluate on the nuScenes validation split. Following common practices, we employ Fréchet Inception Distance (FID)~\cite{fid} and Fréchet Video Distance~\cite{fvd} to quantify the quality of generated scenes. Additionally, we conduct a quantitative analysis to evaluate the model's sensitivity to various motion controls.

\subsection{Implementation Details}
We initialize the 3D VAE and DiT from CogVideoX pretrained weights~\cite{cogvideox}. We construct the trajectory vocabulary using K-means clustering and set the number of trajectory anchors to 256, following WoTE~\cite{li2025wote}. The visual adapter and multi-modal trajectory encoder in the planner are inherited from the pretrained TA-DWM. In Phase 1, the TA-DWM is trained on a combined dataset of driving videos from nuPlan~\cite{caesar2021nuplan} and nuScenes~\cite{nuscenes}. Optimization is performed on 16 NVIDIA A100 GPUs with a batch size of 32. In Phase 2, we freeze the encoders and train the planner for 50 epochs on the NAVSIM navtrain split using 8 NVIDIA 3090 GPUs with a batch size of 256. To train FAR, we freeze the trajectory planner and optimize FAR for 10 epochs using PDMS as the oracle for preference supervision. At inference time, WorldDrive does not perform explicit future scene generation, enabling real-time planning. Additional implementation details are provided in the supplementary material.

\begin{table}[!ht]
    \centering
    \setlength{\tabcolsep}{3pt}
    \resizebox{0.47\textwidth}{!}{
    \begin{tabular}{ccc|cccc|c}
    \toprule[1.5pt]
    VAE & TA-DWM  & TA-DWM   & \multirow{2}{*}{NC} & \multirow{2}{*}{DAC} & \multirow{2}{*}{TTC} & \multirow{2}{*}{EP}  & \multirow{2}{*}{PDMS}  \\ 
    Pretrain & Vision & Motion  &   & & &  & \\ 
    \midrule
    $\times$ & $\times$ & $\times$ & 69.0 & 57.6 & 57.2  & 28.7 & 31.4 \\
    $\checkmark$ & $\times$ & $\times$ &  97.6 & 93.6 & 93.2  & 79.5 & 84.9 \\
    $\checkmark$ & $\checkmark$ & $\times$ &  97.9 & 94.1 & 93.3  & 80.7  & 85.8 \\
    $\checkmark$ & $\checkmark$ & $\checkmark$ &  98.4 & 95.1 & 94.7  & 80.4  &  86.9 \\
    \bottomrule[1.5pt]
\end{tabular}}
\caption{\textbf{Planner with different pretrained representation}. ``VAE Pretrain'' means initializing with the 3D VAE weights from CogVideoX. ``TA-DWM Vision'' and ``TA-DWM Motion'' mean initializing the vision and motion encoders from TA-DWM.}
\label{tab:rep}
\end{table}


% ----------------------- traj part -----------------------
\subsection{WorldDrive on Trajectory Planning}
\noindent\textbf{Performance comparison with SOTA methods.}~Table~\ref{tab:closeloop_turn} compares WorldDrive with state-of-the-art methods on the NAVSIM navtest split. WorldDrive achieves the best PDMS among vision-only methods. Despite relying on a single-view camera, it achieves a PDMS of 88.1, surpassing the previous best single-view method ImagiDrive~\cite{li2025imagidrive} and Epona~\cite{epona} and also outperforms leading multi-view approaches such as FSDrive~\cite{zeng2025fsdrive} and World4Drive~\cite{world4drive}. Notably, our single-view vision-only framework is competitive with multi-modal methods such as DiffusionDrive~\cite{diffusiondrive} and WoTE~\cite{li2025wote}.
When trained on the full navtrain split, WorldDrive further improves PDMS to 89.0 and outperforms the strong world model baseline PWM~\cite{zhao2025pwm} and DriveVLA-W0~\cite{drivevlaw0}. In the best-of-N setting, we use the oracle scorer to select the best trajectory from six candidates to investigate the upper bound, and WorldDrive achieves 93.6 PDMS. This oracle result is reported only to probe the upper bound and to verify that WorldDrive generates a multi-modal candidate set, and it is not used for inference.

Table~\ref{table:navhar} reports results on NAVSIM-v2 navhard split using EPDMS and its sub-metrics under reactive traffic and pseudo closed-loop simulation. WorldDrive achieves the best EPDMS using only a single-view camera, outperforming multi-view baselines such as LTF and DiffusionDrive. The gains are consistent across most compliance and safety-related sub-metrics. We observe that some comfort-related metrics remain challenging, suggesting further effort for improving long-horizon efficiency while preserving safety.

As a supplementary evaluation, we present a comparison on nuScenes in Table~\ref{tab:openloop_turn}. Within the world model-based methods, WorldDrive demonstrates highly competitive results. These results corroborate the effectiveness of our framework, confirming that the synergistic optimization of visual and motion representations within the TA-DWM translates into robust planning capabilities.

\noindent\textbf{Impact of Pre-training Strategy.}~Table~\ref{tab:rep} dissects the benefits of the pre-training curriculum in WorldDrive. Training the planner without any pre-trained representation leads to poor driving performance. Initializing with the 3D Causal VAE provides a foundational leap to 84.9 PDMS, highlighting the importance of strong visual priors from large-scale video pretraining. Crucially, inheriting the TA-DWM vision and motion representation yields a further improvement of 0.9 and 1.1 PDMS. These step-wise gains provide evidence that TA-DWM pretraining aligns the vision and motion feature space with downstream planning requirements, improving both overall PDMS and most sub-metrics.

\noindent\textbf{Trajectory Rewarder Strategy Analysis.}~A central challenge in multi-modal planning is selecting the best trajectory from diverse candidates. Table~\ref{tab:rewarder} evaluates how different inputs contribute to the reward mechanism. Simply relying on trajectory features from planner yields a marginal gain, and while efficiency improves, safety metrics slightly degrade. Incorporating FAR with distilled future latents substantially improves overall performance while simultaneously increasing both safety and efficiency. 


\begin{table}[t]
    \centering
    \setlength{\tabcolsep}{6pt}
    \resizebox{0.47\textwidth}{!}{
    \begin{tabular}{cc|cccc|c}
    \toprule[1.5pt]
    Traj feat. & Future feat.  & NC & DAC & TTC & EP & PDMS  \\ 
    \midrule
    $\times$ & $\times$  & 98.4 & 95.1 & 94.7  & 80.4 & 86.9 \\
    $\checkmark$& $\times$  & 98.3 & 95.3 & 94.3  & 81.2 & 87.0 \\
    $\checkmark$ & $\checkmark$  & 98.4 & 96.2 & 95.1 & 81.9 & 88.1 \\
    \bottomrule[1.5pt]
\end{tabular}}
\caption{\textbf{Future-aware rewarder design}. ``Traj feat.'' means using candidate trajectory features from the planner. ``Future feat.'' means using future scene query and future latent distillation.}
\label{tab:rewarder}
\end{table}

\begin{table}[t]
    \centering
    \setlength{\tabcolsep}{6pt}
    \resizebox{0.47\textwidth}{!}{
    \begin{tabular}{c|c c c c c|c}
    \toprule[1.5pt]
    Method & NC & DAC & TTC & EP & PDMS & Latency\\
    \midrule
    PWM$\dagger$ & 98.0 & 95.1 & 94.1 & 82.4 & 87.3 & 570ms \\
    PWM & 98.6 & 95.9 & 95.4 & 81.8 & 88.1 & 850ms \\
    WorldDrive & 98.4 & 96.8 & 95.2  & 83.3 & 89.0 & 53ms \\ 
    \bottomrule[1.5pt]
\end{tabular}}
\caption{\textbf{Inference latency comparison}. Latency is measured on a single NVIDIA A800 GPU. ``$\dagger$'' means without future frame forecast.}
\label{tab:latency}
\end{table}

\noindent\textbf{Latency Analysis.}~Table~\ref{tab:latency} reports inference efficiency. DWM-based planners face a practical trade-off: incorporating future forecasts improves performance but increases inference latency. WorldDrive avoids this by distilling future latents during training and eliminating explicit future scene generation at inference. WorldDrive achieves strong planning performance with low latency, meeting real-time requirements while retaining benefits of predictive foresight.

\begin{figure}[t]
    \centering
    \includegraphics[width=1.0\linewidth]{imgs/similarity.pdf}
    \caption{\textbf{Quantitative Analysis of Motion Sensitivity.}~The similarity between scene representations is inversely correlated with the geometric distance. The sensitivity to both large~(a) and small~(b) deviations is amplified with further training.}
    \label{fig:similarity}
\end{figure}

\begin{table*}[t]
    \centering
    
    \resizebox{0.98\textwidth}{!}{
    \begin{tabular}{c|cccccc |c}
    \toprule[1.5pt]
    Metric       & DriveDreamer~\cite{drivedreamer} & WoVoGen~\cite{wovogen} & GenAD~\cite{genad}   & Drive-WM~\cite{drivewm} & Vista~\cite{vista} & Driverse~\cite{driverse} & WorldDrive \\
    \midrule
    Resolution     &  128$\times$192 &  256$\times$448 &  256$\times$448 &  192$\times$384 &  480$\times$832 &  480$\times$832 & 256$\times$512 \\
    \midrule
    FID            &     52.6     &   27.6   &   15.4   &  15.2  &   18.2   &   20.1   & 12.8  \\
    FVD            &     452.0    &   417.7  &   184.0  &  122.7 &   158.0  &   143.5  & 131.7 \\
    \bottomrule[1.5pt]
    \end{tabular}}
    \caption{\textbf{nuScenes validation split generated videos comparison}.}
    \label{fig:gen}
\end{table*}

\begin{figure*}[t]
    \centering
    \includegraphics[width=0.98\linewidth]{imgs/planning_result_v9.pdf}
    \caption{\textbf{Qualitative planning result of WorldDrive on NAVSIM navtest split}. (a)~Planning result and the corresponding generated future scene with different trajectories. (b)~Top-10 Multi-modal planning trajectories.}
    \label{fig:planning}
\end{figure*}

\begin{figure*}
    \centering
    \includegraphics[width=0.98\linewidth]{imgs/simulation_result_v8.pdf}
    \caption{\textbf{Qualitative simulation result on nuScenes}. Scenes are generated with three trajectory conditions: the expert trajectory (green), and non-expert alternatives. The results show a tight coupling between the input motion and the generated scenes.}
    \label{fig:visual}
\end{figure*}

\subsection{WorldDrive on Scene Generation}
\noindent\textbf{Generation Quality on nuScenes.}~Table~\ref{fig:gen} presents quantitative results on the nuScenes validation split. WorldDrive attains this performance without relying on the massive training data~\cite{vista} or additional motion-alignment modules~\cite{driverse}, suggesting that robust VAE priors combined with structured trajectory conditioning can yield high-fidelity simulations with a relatively simple design.

\noindent\textbf{Motion sensitivity.}~To quantify the controllability of TA-DWM, we analyze the correlation between motion deviations and the variations of the generated latents. We compute the cosine similarity between latents conditioned on the expert trajectory and those conditioned on trajectory anchors. Fig.~\ref{fig:similarity} shows a clear inverse correlation where latent similarity decreases as the geometric distance increases. Comparing the curves at 10k and 100k iterations reveals that longer training amplifies this sensitivity. We also study the impact of the multi-modal trajectory encoder. Top-5 conditioning consistently yields stronger discrimination than the Top-1 baseline, with the most pronounced improvement in the small-deviation regime~(Fig.~\ref{fig:similarity}(b)).


\subsection{Qualitative Results}
\textbf{Qualitative Results of Planning.}~Fig.~\ref{fig:planning} visualizes WorldDrive’s planning behavior. The planner produces physically feasible trajectories, but the base policy can deviate in challenging cases. With FAR enabled, WorldDrive re-scores candidates using distilled future information and selects a safer trajectory. As shown in the visual simulation panels, TA-DWM can be selectively invoked to generate future scenes for case analysis. Fig.~\ref{fig:planning}~(b) further shows the multi-modal trajectory predictions, providing qualitative evidence that the unified representation supports diverse yet feasible planning behaviors.

\noindent\textbf{Qualitative Results of Scene Generation.}~We qualitatively evaluate TA-DWM by visualizing future scenes generated under different motion conditions in Fig.~\ref{fig:visual}. We condition on the expert trajectory and two counterfactual trajectories. The results show that the generated futures closely follow the specified motion, producing plausible sequences even for counterfactual intentions such as heading toward a non-drivable area or approaching a key traffic participant.







